% Extra preamble for the Pandoc-generated report (see tools/build_report.sh).
% Keep this file free of document-specific text: it only fixes layout that the
% default Pandoc template does not handle well for wide tables and code.

% Latin text and code need a font with broad Unicode coverage: the report
% contains >=, <=, approximately, Greek letters and box-drawing characters
% (tree diagrams) in code blocks, which Latin Modern lacks. CJK is still
% handled by ctex/Fandol; these fonts only cover the Latin/monospace side.
\usepackage{fontspec}
\setmainfont{DejaVu Sans}
\setsansfont{DejaVu Sans}
\setmonofont{DejaVu Sans Mono}[Scale=0.9]

\usepackage{xcolor}
\usepackage{graphicx}
\setkeys{Gin}{width=\linewidth,keepaspectratio}

% Tables: the report has several wide comparison tables. `longtable` handles
% page breaks; \small + a slightly larger row height keeps them readable.
\usepackage{booktabs}
\usepackage{longtable}
\usepackage{array}
\usepackage{ragged2e}
\usepackage{etoolbox}
\AtBeginEnvironment{longtable}{\scriptsize}
\setlength{\tabcolsep}{4pt}
\renewcommand{\arraystretch}{1.15}

% Code blocks: Pandoc's default writer uses a `Highlighting` fancyvrb
% environment. Redefine it at begin-document (after Pandoc's macro) to break
% long lines, so shell commands do not run off the page. Using the default
% writer also keeps inline code as \texttt{...}, which preserves the spaces
% in `mse count` (the deprecated --listings path collapses them).
\usepackage{fvextra}
\AtBeginDocument{%
  \DefineVerbatimEnvironment{Highlighting}{Verbatim}{breaklines,commandchars=\\\{\}}%
}

% The TOC spans the first pages; ctex keeps a "目录" running header alive
% afterwards. Plain pages avoid a misleading header on the body pages.
\pagestyle{plain}

% Links and PDF metadata.
\usepackage{hyperref}
\hypersetup{
  colorlinks=true,
  linkcolor=blue!55!black,
  urlcolor=blue!55!black,
  citecolor=blue!55!black,
  bookmarksnumbered=true,
}
